% Generated draft: requires mathematical and source review.
\documentclass[11pt]{article}
\usepackage[margin=25mm]{geometry}
\usepackage{fontspec}
\setmainfont{Latin Modern Roman}
\usepackage{amsmath,amssymb,mathtools,unicode-math}
\setmathfont{Latin Modern Math}
\usepackage{enumitem,needspace,xurl,hyperref}
\setlength{\parindent}{0pt}
\setlength{\parskip}{5pt}
\setlength{\emergencystretch}{3em}
\newcommand{\N}{\mathbb{N}}
\newcommand{\Z}{\mathbb{Z}}
\newcommand{\Q}{\mathbb{Q}}
\newcommand{\R}{\mathbb{R}}
\newcommand{\C}{\mathbb{C}}
\newcommand{\F}{\mathbb{F}}
\newcommand{\A}{\mathbb{A}}
\newcommand{\PP}{\mathbb{P}}
\newcommand{\E}{\mathbb{E}}
\newcommand{\eps}{\varepsilon}
\newcommand{\dd}{\,\mathrm{d}}
\begin{document}
\section*{UnsolvedMath ID: 30003791}
\textbf{Source problem code:} catalogue entry
\subsection{Definitions and mathematical statement}
Let \(p\) be a prime and let \(F\) be a field of characteristic \(p\). For the standard interpretation of Kato--Milne cohomology, write \(\Omega_F^r=\bigwedge_F^r\Omega^1_{F/\mathbb F_p}\), with \(\Omega_F^0=F\), and let \(\mathrm d\) denote the exterior derivative. For \(n\geq 1\), set
\[
H_p^n(F)=\operatorname{coker}\!\left(\wp:\Omega_F^{n-1}\longrightarrow \Omega_F^{n-1}/\mathrm d\Omega_F^{n-2}\right),
\]
where \(\Omega_F^{-1}=0\) and the additive map \(\wp\) is specified on logarithmic differential expressions by
\[
\wp\!\left(a\,\frac{\mathrm d b_1}{b_1}\wedge\cdots\wedge\frac{\mathrm d b_{n-1}}{b_{n-1}}\right)
= (a^p-a)\,\frac{\mathrm d b_1}{b_1}\wedge\cdots\wedge\frac{\mathrm d b_{n-1}}{b_{n-1}}
\pmod{\mathrm d\Omega_F^{n-2}},
\]
for \(a\in F\) and \(b_i\in F^\times\). The classes of these logarithmic differential expressions are the generating differential symbols. For \(n=1\), the empty wedge is \(1\), so \(H_p^1(F)=F/\{a^p-a:a\in F\}\).

The supplied record asks whether the following implications hold for every field \(F\) of characteristic \(p\) and every degree \(n\) in the intended range:
\[
\begin{aligned}
H_p^n(F)\text{ is linked}&\quad\Longrightarrow\quad H_p^{n+2}(F)=0,\\
H_p^n(F)\text{ is triple linked}&\quad\Longrightarrow\quad H_p^{n+1}(F)=0.
\end{aligned}
\]
Here vanishing means that the indicated cohomology group is the zero group. These are two separate questions.

\textbf{Unresolved formulation qualifications.} The supplied record says that linkage and triple linkage are defined through generating differential symbols, but it does not supply either definition. In particular, it does not specify the common factors, permitted symbol presentations, or quantifiers required for linkage. Accordingly, the two displayed implications retain the source's predicates and are schematic targets pending recovery of those definitions; no particular linkage convention is imposed here. The record also does not specify the intended range of \(n\), including whether degree \(1\) is included. The standard cohomology notation above does not resolve that omission.
\subsection{Short English statement}
The problem asks whether linkage of generating symbols in degree \(n\) forces Kato--Milne cohomology to vanish in degree \(n+2\), and whether triple linkage forces vanishing already in degree \(n+1\). The precise linkage definitions and intended degree range are absent from the supplied record, so a fully specified formulation requires human review.
\subsection{Sources}
\begin{itemize}
\item Local supplied UnsolvedMath registry record, permanent ID \(30003791\), source problem code \texttt{catalogue entry}, in the collection \emph{Oberwolfach Reports Open Problems}. No individual proposer is identified; the supplied proposed year is 2018.
\item Supplied report citation: \emph{Quadratic Forms and Related Structures over Fields (2019)}, Oberwolfach Reports, pp.\ 1235--1291, DOI \texttt{10.4171/owr/2018/21}. \url{https://doi.org/10.4171/owr/2018/21}. The record describes these reports as documenting mathematical workshops and their open-problem discussions.
\item Supplied literature source, labelled ``triple linkage'': \url{https://arxiv.org/abs/1706.04929}.
\item Supplied literature source, labelled ``types of linkage'': \url{https://arxiv.org/abs/1806.03603}.
\item Supplied literature source, labelled ``workshop report'': \url{https://ems.press/content/serial-article-files/46742}.
\item Supplied dated literature curation, checked \texttt{2026-08-21}: the record labels the problem partially solved and states that the \(p=2\) quadratic-Pfister analogues, including triple linkage, are proved, whereas the stated general characteristic-\(p\) Kato--Milne linkage implications remain unresolved. This assessment is supplied dated curation derived from the OWR material, not an independently verified literature assessment; its time-sensitive status and cited claims require independent verification.
\item The registry labels the statement \texttt{corrected\_verified} and reports that it was checked against the official Oberwolfach report. This is a supplied verification claim, not an independent source check performed for this draft.
\end{itemize}
\end{document}
